% =====================================================================
% CORRECCIÓN MÍNIMA DE UNA AFIRMACIÓN — frente de archivos
% =====================================================================
% Se mantiene NapierOne como conjunto de datos. Lo único que se corrige
% es que el texto no afirme algo que no se verificó.
%
% DÓNDE: `resultados.tex`, subsección de limitación del Experimento 2c.
% (Archivo aparte porque esa parte la está editando otra persona.)
%
% DICE HOY:
%   «...y su evaluación requeriría un conjunto de datos con múltiples
%    campañas por familia, no disponible en la actualidad.»
%
% PROBLEMA: «no disponible en la actualidad» es una afirmación sobre el
% mundo, y no se comprobó que sea cierta. Existen colecciones públicas de
% material cifrado cuyo contenido este trabajo no revisó.
%
% CAMBIO PROPUESTO (dos palabras): reemplazar la coma final por
%
%   «...y su evaluación requeriría un conjunto de datos con múltiples
%    campañas por familia, del que no se dispuso para este trabajo.»
%
% Con eso la frase dice exactamente lo que se sabe: que no lo tuvimos.
% No afirma que no exista, y no obliga a defender esa afirmación.
% =====================================================================
